% Figure: pressure-enthalpy (p-h) diagram for carbon dioxide running a
% transcritical refrigeration cycle. The dome peaks at the critical point
% (31 C, 7.38 MPa); the high-side pressure (gas cooler) sits ABOVE the
% critical pressure, so heat rejection follows a supercritical isobar that
% never crosses the dome, while the low-side (evaporator) stays subcritical
% inside the dome. The four state points are the worked CO2 example:
% 1 saturated vapour at evaporator pressure (~3.5 MPa, h ~ 435),
% 2 after isentropic compression to the gas-cooler pressure (10 MPa, h ~ 480),
% 3 gas-cooler exit (10 MPa, 35 C, h ~ 250, supercritical liquid-like),
% 4 after isenthalpic throttle back to evaporator pressure (h ~ 250, two-phase).
\begin{figure}[htbp]
\centering
\definecolor{engteal}{RGB}{30,120,120}
\begin{tikzpicture}
\begin{axis}[
  width=12cm, height=8cm,
  xlabel={specific enthalpy $h$ (kJ/kg)},
  ylabel={pressure $p$ (MPa, log scale)},
  ymode=log,
  xmin=150, xmax=520, ymin=2, ymax=14,
  axis lines=left,
  log ticks with fixed point,
  tick label style={font=\scriptsize},
  label style={font=\small},
  clip=true,
]
% saturation dome for CO2, peaking at the critical point (h ~ 332, p = 7.38 MPa).
% left limb (saturated liquid), rising to the critical point:
\addplot[engblue, very thick, smooth] coordinates {
  (180,2.0) (200,3.0) (220,4.2) (250,5.6) (290,6.8) (332,7.38)
};
\addplot[only marks, mark=*, engblack, mark size=1.8pt] coordinates {(332,7.38)};
% right limb (saturated vapour), descending from the critical point:
\addplot[engblue, very thick, smooth] coordinates {
  (332,7.38) (380,6.8) (415,5.6) (435,4.2) (445,3.0) (450,2.0)
};
\node[engblack, font=\scriptsize, anchor=south] at (axis cs:332,7.6) {critical point};
\node[engblue, font=\scriptsize] at (axis cs:310,3.2) {two-phase};
\node[enggray, font=\scriptsize, anchor=west] at (axis cs:175,11.5) {supercritical};
% critical-pressure reference line (dashed) so the reader sees the high side is above it
\addplot[enggray, dashed, thick] coordinates {(150,7.38) (520,7.38)};
\node[enggray, font=\tiny, anchor=west] at (axis cs:152,7.9) {$p_c=7.38$ MPa};
% transcritical cycle state points:
% 1: saturated vapour, evaporator p = 3.5 MPa, h1 = 435
% 2: after isentropic compression to 10 MPa, h2 = 480 (supercritical)
% 3: gas-cooler exit, 10 MPa, 35 C, h3 = 250 (supercritical, dense)
% 4: after throttle to 3.5 MPa, h4 = 250, two-phase
\addplot[only marks, mark=*, engred, mark size=2.0pt] coordinates {
  (435,3.5) (480,10) (250,10) (250,3.5)
};
% 1 -> 2 compression
\addplot[engred, very thick] coordinates {(435,3.5) (480,10)};
% 2 -> 3 gas cooling (constant supercritical pressure, leftwards)
\addplot[engorange, very thick] coordinates {(480,10) (250,10)};
% 3 -> 4 throttle (constant h, downwards across the dome)
\addplot[engred, very thick] coordinates {(250,10) (250,3.5)};
% 4 -> 1 evaporation (constant p, rightwards inside the dome)
\addplot[engteal, very thick] coordinates {(250,3.5) (435,3.5)};
\node[engblack, font=\scriptsize, anchor=north] at (axis cs:435,3.35) {1};
\node[engblack, font=\scriptsize, anchor=west]  at (axis cs:483,10) {2};
\node[engblack, font=\scriptsize, anchor=east]  at (axis cs:247,10) {3};
\node[engblack, font=\scriptsize, anchor=north] at (axis cs:250,3.35) {4};
\node[engred, font=\scriptsize, anchor=west] at (axis cs:455,6.0) {compressor};
\node[engorange, font=\scriptsize, anchor=south] at (axis cs:360,10.2) {gas cooler};
\node[engteal, font=\scriptsize, anchor=north] at (axis cs:340,3.35) {evaporator};
\end{axis}
\end{tikzpicture}
\caption[Pressure-enthalpy diagram of a CO\textsubscript{2} transcritical
cycle]{The pressure-enthalpy diagram of carbon dioxide carrying a
transcritical refrigeration cycle. The saturation dome (blue) peaks at the
low critical point of CO\textsubscript{2} ($31\,\degC$, $7.38\,\text{MPa}$).
The evaporator (state~4 to state~1) sits inside the dome at a subcritical
pressure, so the cold side behaves like an ordinary boiling refrigerant.
The compressor lifts the vapour past the critical pressure (dashed line) to
the gas-cooler pressure of $10\,\text{MPa}$, so heat rejection from state~2
to state~3 follows a supercritical isobar that never enters the dome and
never condenses; the fluid cools as a dense supercritical phase rather than
condensing at constant temperature. The throttle (state~3 to state~4) drops
the dense supercritical fluid back into the two-phase region. The cooling
capacity is the evaporator width $q_L=h_1-h_4$ and the work is the
compressor width $w_c=h_2-h_1$; the gas-cooler exit enthalpy, set by the
high-side pressure and the gas-cooler outlet temperature, controls both,
which is why the high-side pressure is a free variable the designer
optimises rather than a quantity the condenser fixes.}
\label{fig:vol04ch02:co2-transcritical-cycle}
\end{figure}
